\section{Probability Theory}

\begin{definition}
    A \textbf{discrete probability space} is a set of elementary events often denoted as 
    $$\Omega = \{ \omega_1, \omega_2, \dots, \omega_n \}$$
    Every \textbf{elementary event} $\omega_i$ has a probability $P[\omega_i] \in [0,1]$ such that
    $$\sum_{\omega_i \in \Omega} P[\omega_i] = 1$$
\end{definition}

\begin{remark}
    A finite probability space, in which every elementary event has the same probability is called a
    \textbf{Laplace-Space}. Observe that for a Laplace-Space $\Omega$ and an arbitrary event $E$ we have:
    $$P[E] = \frac{\vert E \vert}{\vert \Omega \vert}$$
\end{remark}

\begin{definition}
    An \textbf{Event} $E \subseteq \Omega$ is a set of elementary events. The probability of an event is defined as:
    $$P[E] = \sum_{\omega_i \in E} P[\omega_i]$$
    The complementary event $\bar{E} = \Omega \setminus E$ has probability:
    $$P[\bar{E}] = 1 - P[E]$$
\end{definition}

\begin{theorem}
    For two events $A, B$ in a probability space $\Omega$ we have:
    \begin{itemize}
        \item $P[\emptyset] = 0$ and $P[\Omega] = 1$
        \item $P[A] \in [0,1]$
        \item $P[\bar{A}] = 1 - P[A]$
        \item If $A \subset B$ then $P[A] \le P[B]$
    \end{itemize}
\end{theorem}

\begin{theorem}
    For events $A_1, A_2, ... A_n$ it holds
    $$P[A_1 \cup A_2 \cup \dots \cup A_n] \leq \sum_{i=1}^n P[A_i]$$
    with equality iff $A_1, A_2, ... A_n$ are pairwise disjoint.
\end{theorem}

\begin{proof}
    For $i \in \{1, \dots, n\}$ let $B_i = A_i \setminus \bigcup_{j=1}^{i-1} A_j$. Then $P[B_i] \leq P[A_i]$ and for $i \neq j$ we have $B_i \cap B_j = \emptyset$.
    $$P[\bigcup_{i=1}^n A_i] = P[\bigcup_{i=1}^n B_i] = \sum_{i=1}^n P[B_i] \leq \sum_{i=1}^n P[A_i]$$
\end{proof}

\begin{theorem}
    \textbf{Inclusion-Exclusion Principle}: For events $A_1, A_2, ... A_n$ with $n \geq 2$ it holds
    $$P[\bigcup_{i=1}^n A_i] = \sum_{i=1}^n (-1)^{i+1} P[\bigcap_{j=1}^i A_j]$$
\end{theorem}

\begin{proof}
    To proof this theorem we look at an arbitrary $\omega \in \bigcup_{i=1}^n A_i$. If $\omega$ is in exactly
    $k$ of the events $A_i$, then it is counted $N_k$ times where
    $$N_k = k - \binom{k}{2} + \binom{k}{3} - \dots + (-1)^{k+1} \binom{k}{k}$$
    We show that $N_k = 1$ for all $k \in \{1, 2, \dots, n\}$. For that we use the binomial theorem:
    $$(x-y)^k = \sum_{i=0}^k \binom{k}{i} x^{k-i} (-y)^i$$
    Now let $x=1$ and $y=1$:
    $$(1-1)^k = \sum_{i=0}^k \binom{k}{i} (-1)^i = \binom{k}{0} - \binom{k}{1} + \binom{k}{2} - \dots + (-1)^k
    \binom{k}{k} = 1- N_k$$
    Therfore:
    $$N_k = 1 - (1-1)^k = 1$$
\end{proof}

\begin{example}
    We shuffle a deck of cards and want to know the probability that any of the cards is at the same place as it was before.
    Let $A_i$ be the event that the $i$-th card is at the same place as it was before. Then we want to find
    $$P[A_1 \cup A_2 \cup \dots \cup A_{52}]$$
    We define $A_S$ as the event that all cards $i \in S \subseteq[n]$ are at the same place as they were before.
    $$P[A_S] = \frac{(n - \vert S \vert)!}{n!}$$
    Using the inclusion-exclusion principle we get
    $$P[\bigcup_{i=1}^n A_i] = \sum_{i=1}^n A_i - \sum_{S \subset [n], \vert S \vert = 2} A_S + \sum_{S \subset [n], \vert S \vert = 3} A_S - \dots$$
    $$= n \frac{(n-1)!}{n!} - \binom{n}{2} \frac{(n-2)!}{n!} + \binom{n}{3} \frac{(n-3)!}{n!} - \dots$$
    $$= 1 - \frac{1}{2!} + \frac{1}{3!} - \dots + (-1)^{n+1} \frac{1}{n!} \approx 1 - \frac{1}{e} \approx 0.632$$
\end{example}

\subsection{Conditional Probability}

\begin{definition}
    Let $A$ and $B$ be events with $P[B] > 0$. The \textbf{conditional probability} of $A$ given $B$ is defined as:
    $$P[A \mid B] = \frac{P[A \cap B]}{P[B]}$$
\end{definition}

\begin{remark}
    Let $A \subseteq \Omega$ be an event, then the following always hold
    \begin{itemize}
        \item $P[A \mid \Omega] = P[A]$
        \item $P[A \mid A] = 1$
        \item $P[A \mid \bar{A}] = 0$
        \item $P[A \cap B] = P[A \mid B] P[B] = P[B \mid A] P[A]$
    \end{itemize}
\end{remark}

\begin{example}
    \textbf{Two-Siblings Problem (Zweikinderproblem)}:
    Assume a family has two children, and both genders are equally likely at birth. We define the sample space as:
    $$\Omega = \{ \text{gg}, \text{gb}, \text{bg}, \text{bb} \}$$
    where the letters represent the gender (g for girl, b for boy) in order of birth. Each elementary event has probability $1/4$.
    We want to find the probability that both children are girls ($A = \{ \text{gg} \}$).
    
    \begin{enumerate}
        \item \textbf{Scenario 1}: We know that the family has at least one girl ($B_{\text{at least one}} = \{ \text{gg}, \text{gb}, \text{bg} \}$). Then:
        $$P[A \mid B_{\text{at least one}}] = \frac{P[A \cap B_{\text{at least one}}]}{P[B_{\text{at least one}}]} = \frac{1/4}{3/4} = \frac{1}{3}$$
        \item \textbf{Scenario 2}: We know that the older child is a girl ($B_{\text{older}} = \{ \text{gg}, \text{gb} \}$). Then:
        $$P[A \mid B_{\text{older}}] = \frac{P[A \cap B_{\text{older}}]}{P[B_{\text{older}}]} = \frac{1/4}{2/4} = \frac{1}{2}$$
    \end{enumerate}
    This shows how critical the exact conditioning information is, as the gender of the older sibling is independent of the younger one, whereas the global statement of "at least one girl" couples the two events.
\end{example}

\begin{theorem}
    \textbf{Multiplication-Theorem: } Let $A_1, A_2, .. A_n$ be events. If $P[A_1 \cap A_2 \cap \dots \cap A_{n-1}] > 0$ then it holds:
    $$P[A_1 \cap A_2 \cap \dots \cap A_n] = P[A_1] \cdot P[A_2 \mid A_1] \cdot P[A_3 \mid A_1 \cap A_2] \dots P[A_n \mid A_1 \cap A_2 \cap \dots \cap A_{n-1}]$$
    
\end{theorem}

\begin{proof}
    We can prove the correctness of this theorem by using the definition of conditional probability:
    $$P[A_1] \cdot P[A_2 \mid A_1] \cdot P[A_3 \mid A_1 \cap A_2] \dots P[A_n \mid A_1 \cap A_2 \cap \dots \cap A_{n-1}]$$
    $$ = P[A_1] \cdot \frac{P[A_1 \cap A_2]}{P[A_1]} \cdot \frac{P[A_1 \cap A_2 \cap A_3]}{P[A_1 \cap A_2]} \dots \frac{P[A_1 \cap A_2 \cap \dots \cap A_n]}{P[A_1 \cap A_2 \cap \dots \cap A_{n-1}]}$$
    $$ = P[A_1 \cap A_2 \cap \dots \cap A_n]$$
\end{proof}

\begin{example}
    \textbf{Birthday Problem: } What is the probability of two people in a room of $m$ having the same birthday? We find the probability that none of them have the same birthday this is 
    $$\frac{364}{365} \cdot \frac{363}{365} \cdot \dots \cdot \frac{365-m+1}{365}$$
\end{example}

\begin{theorem}
    \textbf{Theorem of total probability: }Let $A_1, ... A_n$ be pairwise disjoint events and let $B \subseteq A_1 \cup ... \cup A_n$ then it holds:
    $$P[B] = \sum_{i=1}^n P[B \mid A_i] P[A_i]$$
\end{theorem}

\begin{example}
    \textbf{Goat Problem: } Assume three doors of which we know that one hides a sportscar and the other two a goat. You pick one arbitrary door, then one of the other two opens and you see a goat. Should you now switch to the other door that is left or should you stay with the one you are?
    Let $A_1$ be the event that you picked the door with the car and $A_2$ the event that you are infront of the door with a goat. Further let $B$ be the event that you win if you change doors. Now observe
    $$Pr[B \mid A_1] = 0$$
    $$Pr[B \mid A_2] = 1$$
    $$Pr[A_1] = 1/3$$
    $$Pr[A_2] = 2/3$$
    Using the theorem above we get:
    $$Pr[B] = Pr[B \mid A_1] Pr[A_1] + Pr[B \mid A_2] Pr[A_2] = 0 \cdot 1/3 + 1 \cdot 2/3 = 2/3$$
\end{example}

\begin{theorem}
    \textbf{Bayes' Theorem: } Let $A_1, A_2 ... A_n$ be pariwise disjoint events and $B \subseteq A_1 \cup A_2 ... \cup A_n$ with $P[B] > 0$. Then:
    $$P[A_i \mid B] = \frac{P[B \mid A_i] P[A_i]}{\sum_{j=1}^n P[B \mid A_j] P[A_j]}$$
\end{theorem}

\begin{proof}
    To prove bayes theorem we use the theorem of total probability:
    $$P[A_i \mid B] = \frac{P[B \cap A_i]}{P[B]} = \frac{P[B \mid A_i] P[A_i]}{\sum_{j=1}^n P[B \mid A_j] P[A_j]}$$
\end{proof}

    \begin{example}

\textbf{Desease Test: }We test a person for a desease that has a probability of $P[A] = 0.01$ to occur.
    Let $A$ be the event that the tested person has the desease. Further let $B$ be the event that the test
    was positive. Assume we know the sensitivity of the test $P[B \mid A] = 0.72$ and the specificity
    $P[\bar{B} \mid \bar{A}] = 0.98$. We want to find the probability that a person has the desease under
    the condition that the test is positive:
    $$P[A \mid B] = \frac{P[B \mid A] P[A]}{P[B \mid A] P[A] + P[B \mid
    \bar{A}] P[\bar{A}]} = \frac{0.72 \cdot 0.01}{0.72 \cdot 0.01 + 0.02 \cdot 0.99} = \frac{0.0072}{0.027}
    \approx 0.267$$
    Here we see that even if the test is quite accurate, the probability that a person has the desease
    is still quite low. This is because the probability of the desease is very low. If the probability of
    the desease was higher, the probability of the test being positive would be higher.
    
\end{example}

\subsection{Independence}

\begin{definition}
    Two events $A$ and $B$ are \textbf{independent} if $P[A \cap B] = P[A] P[B]$ or equivalently $P[A \mid B] = P[A]$.
\end{definition}

\begin{definition}
    Events $A_1, A_2, \dots, A_n$ are independent if
    $$\forall I \subseteq \{1, 2, \dots, n\}\quad P[\bigcap_{i \in I} A_i] = \prod_{i \in I} P[A_i]$$
\end{definition}

\begin{remark}
    For infinite families of events it is not enough to check the condition for all finite subsets.
\end{remark}

\begin{lemma}
    Let $A_1, A_2 \dots A_n$ be events and let $A_i^{0} \coloneqq \bar{A_i}$ and $A_i^{1} \coloneqq A_i$ for $i \in \{1, 2, \dots, n\}$.
    Then $A_1, A_2 \dots A_n$ are independent iff for all $(s_1, \dots s_n) \in \{0,1\}^{n}$ it holds that:
    $$P[\bigcap_{i=1}^n A_i^{s_i}] = \prod_{i=1}^n P[A_i^{s_i}]$$
\end{lemma}

\begin{proof}
    $\impliedby: \quad$ We proof this direction by induction over the number of events with exponents $s_i = 0$. For the basecase $s_1 = s_2 \dots = 1$ it holds trivially.
    $$P[\bar{A_1} \cap A_2^{s_2} \cap A_3^{s_3} \dots] = P[A_2^{s_2} \cap A_3^{s_3} \dots] - P[A_1 \cap A_2^{s_2} \cap A_3^{s_3} \dots]$$
    Using the incuction hypothesis we get:
    $$= \prod_{i=2}^n P[A_i^{s_i}] - P[A_1]\prod_{i=2}^n P[A_i^{s_i}] = (1-P[A_1])\prod_{i=2}^n P[A_i^{s_i}] = P[\bar{A_1}]\prod_{i=2}^n P[A_i^{s_i}]$$
    $\implies: \quad$ ...
\end{proof}

\begin{lemma}
    Let $A, B$ and $C$ be independent events then so are $A \cap B$ and $C$ and also $A \cup B$ and $C$.
\end{lemma}

\subsection{Random Variables}

\begin{definition}
    A \textbf{random variable } is a map $X: \Omega \to \mathbb{R}$ where $\Omega$ is the set of all events of a probability space.
\end{definition}

\begin{example}
    The event that the random variable X takes a value less than or equal to 5 is denoted as:
    $$X \leq 5 \coloneqq \{\omega \in \Omega \mid X(\omega) \leq 5\}$$
\end{example}

\begin{definition}
    The \textbf{probability mass function} of a discrete random variable $X$ is defined as:
    $$f_X(x) = P(X=x)$$
\end{definition}

\begin{definition}
    The \textbf{distribution function} of a discrete random variable $X$ is defined as:
    $$F_X(x) = P(X \leq x)$$
\end{definition}

\subsubsection{Expectation}

\begin{definition}
    Let $X$ be a random variable. The expectation of $X$ is defined as:
    $$
    \mathbb{E}[X] \coloneqq \sum_{\alpha \in W_x} \alpha P(X=\alpha)
    $$
    aslong as the sum is absolutely convergent, else we say the expectation is undefined.
\end{definition}

\begin{lemma}
    For any random variable $X$ we have:
    $$ \mathbb{E}[X] = \sum_{\omega \in \Omega} X(\omega) P(\omega) = \sum_{k=1}^\infty P(X \geq k) - \sum_{k=1}^\infty P(X \leq -k) $$
\end{lemma}

\begin{proof}
    $$\mathbb{E}[X] = \sum_{\alpha \in W_X}\alpha P[X=\alpha] = \sum_{\alpha \in W_X}\alpha\sum_{\omega \in \Omega, X(\omega)=\alpha} P[\omega] = \sum_{\omega \in \Omega} X(\omega) P[\omega]$$
\end{proof}

\begin{theorem}
    Let $X$ be a random variable with values in $W_X \subseteq \mathbb{N}_0$. Then:
    $$\mathbb{E}[X] = \sum_{k=1}^\infty P[X \geq k]$$
\end{theorem}

\begin{proof}
    $$\mathbb{E}[X] = \sum_{k=0}^\infty k P[X=k] = \sum_{k=0}^\infty \sum_{j=1}^k P[X=k] = \sum_{j=1}^\infty \sum_{k=j}^\infty P[X=k] = \sum_{j=1}^\infty P[X \geq j]$$
\end{proof}

\begin{example}
    We throw a coin and we know $P[\text{head}] = p$. We want to know $\mathbb{E}[X]$ with $X \coloneqq \text{number of throws until we get head}$. Using the above theorem we get:
    $$\mathbb{E}[X] = \sum_{k=1}^\infty P[X \geq k] = \sum_{k=1}^\infty (1-p)^{k-1} = \frac{1}{1-(1-p)} = \frac{1}{p}$$
\end{example}

\begin{definition}
    For an event $E \subseteq \Omega$ the corresponding indicator-variable $X_E$ is defined as:
    $$ X_E(\omega) \coloneqq \begin{cases} 1 & \text{if } \omega \in E \\ 0 & \text{if } \omega \notin E \end{cases}$$
\end{definition}

\begin{remark}
    For an indicator variable $X_E$ we have:
    $$\mathbb{E}[X_E] = P[E]$$
\end{remark}

\begin{theorem}
    The linearity of expectation says that for random variables $X_1, \dots, X_n$  and $a_1, \dots a_n, b \in \mathbb{R}$ we have:
    $$X \coloneqq \sum_{i=1}^n a_i X_i + b$$
    $$\mathbb{E}[X] = \sum_{i=1}^n a_i \mathbb{E}[X_i] + b$$
\end{theorem}

\begin{remark}
    We can intuitively describe this as "the expectation of a sum is the sum of the expectations".
\end{remark}

\begin{proof}
    $$\mathbb{E}[X] = \sum_{\omega \in \Omega} X(\omega) P[\omega] = \sum_{\omega \in \Omega} (a_1 X_1(\omega) + \dots + a_n X_n(\omega)) P[\omega]$$
    $$ = a_1 \sum_{\omega \in \Omega} X_1(\omega) P[\omega] + \dots + a_n \sum_{\omega \in \Omega} X_n(\omega) P[\omega] = a_1 \mathbb{E}[X_1] + \dots + a_n \mathbb{E}[X_n]$$
\end{proof}

\begin{example}
    We throw a coin $m$ times. Let $X \coloneqq$ number of partial sequences with threetimes heads. What is the expectation of $X$?
    $$X = X_1 + X_2 + \dots + X_{m-2}$$
    where $X_i$ is the indicator variable for the event that the $i$-th partial sequence is three heads.
    $$\mathbb{E}[X_i] = \frac{1}{8}$$
    By the linearity of the expectation we get:
    $$\mathbb{E}[X] = \sum_{i=1}^{m-2} \mathbb{E}[X_i] = \frac{m-2}{8}$$
\end{example}

\begin{remark}
    Recall Stable sets in a graph is a set of vertices that are not adjacent to each other. We will now use a randomized algorithm to determin a maximal stable set.
\end{remark}

\begin{algorithm}
S[v] = 1/0 with probability p
for all u,v in E:
    if S[u] = 1 and S[v] = 1:
        S[v] = 0
return S
\end{algorithm}

this results in a stable set but what is the expected size of S?

\begin{theorem}
    The algorithm produces a stable set of expected size at least $n*P-m*P^2$
\end{theorem}

\begin{proof}
    Let $X$ be the random variable for the size of the stable set $S$ after the first step. Further let $Y$ be the random variable for the amount of edges we looked at in step 2. Observe
    $$\mid S \mid  \geq X - Y$$
    By the linearity of expectation we get:
    $$\mathbb{E}[\mid S \mid] \geq \mathbb{E}[X] - \mathbb{E}[Y]$$
    We know that $\mathbb{E}[X] = nP$. Further we know that $\mathbb{E}[Y] = mP^2$. Thus we get:
    $$\mathbb{E}[\mid S \mid] \geq nP - mP^2$$
\end{proof}

\subsubsection{Variance}

\begin{definition}
    For a random variable $X$ with $\mathbb{E}[X] = \mu$ the variance is defined as:
    $$Var[X] \coloneqq \mathbb{E}[(X-\mu)^2]$$
    The standard deviation of $X$ is defined as $\sigma \coloneqq \sqrt{Var[X]}$.
\end{definition}

\begin{example}
    Let $P[X=-10^{10}] = 1/2$ and $P[X=10^{10}] = 1/2$. Then $\mathbb{E}[X] = 0$ and $Var[X] = 10^{20}$ and $\sigma = 10^{10}$.
\end{example}

\begin{theorem}
    For a random variable $X$ we have $Var[X] = \mathbb{E}[X^2] - \mathbb{E}[X]^2$
\end{theorem}

\begin{proof}
    $$Var[X] = \mathbb{E}[(X-\mu)^2] = \mathbb{E}[X^2 - 2\mu X + \mu^2] = \mathbb{E}[X^2] - 2\mu \mathbb{E}[X] + \mu^2 = \mathbb{E}[X^2] - \mathbb{E}[X]^2 $$
\end{proof}

\begin{lemma}
    For a random variable $X$ and $a, b \in \mathbb{R}$ we have $Var[aX+b] = a^2 Var[X]$
\end{lemma}

\begin{proof}
    \begin{enumerate}
        \item $Var[X+b] = \mathbb{E}[(X+b-\mathbb{E}[X+b])^2] = \mathbb{E}[(X+b-(\mathbb{E}[X]+b))^2] = \mathbb{E}[(X-\mathbb{E}[X])^2] = Var[X]$
        \item $Var[aX] = \mathbb{E}[(aX-\mathbb{E}[aX])^2] = \mathbb{E}[(aX-a\mathbb{E}[X])^2] = \mathbb{E}[a^2(X-\mathbb{E}[X])^2] = a^2 \mathbb{E}[(X-\mathbb{E}[X])^2] = a^2 Var[X]$
    \end{enumerate}
\end{proof}

\begin{definition}
    For a random variable $X$ we call $\mathbb{E}[X^k]$ the $k$-th moment of $X$ and $\mathbb{E}[(X-\mathbb{E}[X])^k]$ the $k$-th central moment of $X$.
\end{definition}

\begin{remark}
    Observe that the first moment is the expectation itself. Further observe that the second central moment is the variance.
\end{remark}

\subsubsection{Multiple Random Variables}

\begin{example}
    Throw a coin and a dice. Let $X$ be the outcome of the coin and $Y$ be the outcome of the dice. What is the expectation of $X+Y$?
\end{example}

\begin{definition}
    Random variables $X_1, X_2, \dots, X_n$ are called independent iff for all $a_1, \dots, a_n \in W_{X_1} \times \dots \times W_{X_n}$ we have:
    $$P[X_1=a_1, \dots, X_n=a_n] = P[X_1=a_1] \dots P[X_n=a_n]$$
\end{definition}

\begin{lemma}
    Observe that events $A_1, \dots, A_n$ are independent iff their indicator variables $X_{A_1}, \dots, X_{A_n}$ are independent.
\end{lemma}

\begin{proof}
    Note that the event "$I_{A_j} = \alpha_j$" is equivalent to the event "$A_j$" if $\alpha_j = 1$ and to the event $\bar{A_j}$ if $\alpha_j = 0$. Same for $A_j^{\alpha_j}$.
    If now $A_1, \dots, A_n$ are independent then for all $\alpha_1, \dots, \alpha_n \in \{0,1\}$ we have:
    $$P[I_{A_1} = \alpha_1, \dots, I_{A_n} = \alpha_n] = P[I_{A_1} = \alpha_1] \dots P[I_{A_n} = \alpha_n]$$
    Using the above note we get:
    $$P[A_1^{\alpha_1} \cap \dots \cap A_n^{\alpha_n}] = P[A_1^{\alpha_1}] \dots P[A_n^{\alpha_n}]$$
\end{proof}

\begin{lemma}
    Let $X_1, \dots, X_n$ be independent random variables and $S_1, \dots, S_n \subseteq \mathbb{R}$ be arbitrary sets then:
    $$P[X_1 \in S_1, \dots, X_n \in S_n] = P[X_1 \in S_1] \dots P[X_n \in S_n]$$
    (Further this implies that for $X_1, \dots, X_n$ independent random variables and $I \subseteq \{1, \dots, n\}$ we have $(X_i)_{i \in I}$ are independent.)
\end{lemma}

\begin{proof}
    $$P[X_1\in S_1, \dots , X_n\in S_n] = \sum_{\alpha_1 \in S_1} \dots \sum_{\alpha_n \in S_n} P[X_1=\alpha_1, \dots, X_n=\alpha_n]$$
    by indepencence of $X_i$ we get:
    $$= \sum_{\alpha_1 \in S_1} \dots \sum_{\alpha_n \in S_n} P[X_1=\alpha_1] \dots P[X_n=\alpha_n]$$
    $$= (\sum_{\alpha_1 \in S_1} P[X_1=\alpha_1]) \dots (\sum_{\alpha_n \in S_n} P[X_n=\alpha_n]) = P[X_1\in S_1] \dots P[X_n\in S_n]$$
\end{proof}

\begin{theorem}
    Let $f_1, \dots, f_n: \mathbb{R} \to \mathbb{R}$ be functions in $\mathbb{R}$ and $X_1, \dots, X_n$ be independent random variables. Then $f_1(X_1), \dots, f_n(X_n)$ are independent too.
\end{theorem}

\begin{proof}
    Let $S_i = \{\beta \mid f_i(\beta) = \alpha_i\}$. Then we have:
    $$P[f_1(X_1) = \alpha_1, \dots, f_n(X_n) = \alpha_n] = P[X_1 \in S_1, \dots, X_n \in S_n]$$
    by the previous lemma we get:
    $$= P[X_1 \in S_1] \dots P[X_n \in S_n] = P[f_1(X_1) = \alpha_1] \dots P[f_n(X_n) = \alpha_n]$$
\end{proof}
    
\subsubsection{Combined Random Variables}

\begin{theorem}
    For two independent random variable $X, Y$ and $Z = X+Y$ we have:
    $$f_Z(\alpha) = \sum_{\beta \in W_X} f_X(\beta) f_Y(\alpha-\beta)$$
\end{theorem}

\begin{remark}
    This can intuitivaly be seen by looking at every possibility for $X = \beta$: for $X+Y = \alpha$ to hold we need $Y = \alpha - \beta$. Since $X$ and $Y$ are independent we can multiply their probabilities and sum that over all $\beta$
\end{remark}

\begin{theorem}
    For random variables $X_1, \dots , X_n$ we have:
    $$\mathbb{E}[X_1 + \dots + X_n] = \mathbb{E}[X_1] + \dots + \mathbb{E}[X_n]$$
\end{theorem}

\begin{remark}
    \textbf{Calculation Rules} for combined random variables
    \begin{enumerate}
        \item $\mathbb{E}[X+Y] = \mathbb{E}[X] + \mathbb{E}[Y]$
        \item $\mathbb{E}[XY] = \mathbb{E}[X] \mathbb{E}[Y]$ if $X, Y$ are independent
        \item $Var[X+Y] = Var[X] + Var[Y]$ if $X, Y$ are independent
    \end{enumerate}
\end{remark}

\begin{theorem}
    For independent random variablex $X_1, \dots , X_n$ we have:
    $$\mathbb{E}[X_1 \dots X_n] = \mathbb{E}[X_1] \dots \mathbb{E}[X_n]$$
    $$Var[X_1 + \dots + X_n] = Var[X_1] + \dots + Var[X_n]$$
\end{theorem}

\begin{proof}
    We prove for $n = 2$ general case can be proved inductively.
    $$\mathbb{E}[(X+Y)^2] = \mathbb{E}[X^2 + 2XY + Y^2] = \mathbb{E}[X^2] + 2\mathbb{E}[XY] + \mathbb{E}[Y^2]$$
    $$(\mathbb{E}[X+Y])^2 = (\mathbb{E}[X] + \mathbb{E}[Y])^2 = \mathbb{E}[X]^2 + 2\mathbb{E}[X]\mathbb{E}[Y] + \mathbb{E}[Y]^2$$
    By subtracting both equations we get:
    $$\mathbb{E}[(X+Y)^2] - (\mathbb{E}[X+Y])^2 = \mathbb{E}[X^2] - \mathbb{E}[X]^2 + \mathbb{E}[Y^2] - \mathbb{E}[Y]^2$$
    Which is equivalent to $Var[X+Y] = Var[X] + Var[Y]$.
\end{proof}

\begin{important}
    Note that $Var[X \cdot Y] = Var[X] \cdot Var[Y]$ does not hold in general.
\end{important}

\begin{theorem}
    Let $X$ be a random variable. For pariwise disjoint events $A_1, \dots, A_n$ with $A_1 \cup \dots \cup A_n = \Omega$  and $P[A_i] > 0 \forall i$ we have:
    $$\mathbb{E}[X] = \sum_{i=1}^n \mathbb{E}[X \mid A_i] P[A_i]$$
\end{theorem}

\begin{proof}
    $$\mathbb{E}[X] = \sum_{\omega \in \Omega} P[X=\omega] \omega = \sum_{i=1}^n \omega \sum_{\omega \in A_i} P[X=\omega | A_i] P[A_i]$$
    $$ = \sum_{i=1}^n A_i \sum_{\omega \in W_X } \omega P[X=\omega | A_i] = \sum_{i=1}^n \mathbb{E}[X \mid A_i] P[A_i]$$
\end{proof}

\begin{theorem}
    \textbf{The waldsche identity}
    Let $N$ and $X$ be independent random variabeles further let
    $$Z \coloneqq \sum_{i=1}^N X_i$$
    where $X_i$ are independent copies of $X$ then:
    $$\mathbb{E}[Z] = \mathbb{E}[N] \mathbb{E}[X]$$
\end{theorem}

\begin{proof}
    $$\mathbb{E}[Z] = \sum_{k=1}^n \mathbb{E}[Z \mid N=k] P[N=k] = \sum_{k=1}^n \mathbb{E}[X_1 + \dots + X_k] P[N=k]$$
    $$= \sum_{k=1}^n k \mathbb{E}[X] P[N=k] = \mathbb{E}[X] \sum_{k=1}^n k P[N=k] = \mathbb{E}[X] \mathbb{E}[N]$$
\end{proof}

\begin{example}
    Throw a dice where $N$ is the face of the dice we get. Afterwards throw a coin $N$-times and let $Z$ be the number of heads. What is the expectation of $Z$?
    $$\mathbb{E}[Z] = \sum_{k=1}^6 P[N=k] \mathbb{E}[Z \mid N=k] = \sum_{k=1}^6 \frac{1}{6} \frac{k}{2} = \frac{1}{12} \sum_{k=1}^6 k = \frac{1}{12} \frac{6 \cdot 7}{2} = \frac{7}{4}$$
\end{example}

\begin{example}
    Throw a coin until we get heads. Let $N$ be the number of throws. Then throw a coin $N$-times and let $Z$ be the number of heads. What is the expectation of $Z$?
    $$\mathbb{E}[N] = 2 (\text{geometric distribution } p = 1/2)$$
    $$Z = \sum_{i=1}^N X_i \quad \text{where Xi are bernulli distributed with} \quad p = 1/2 \text{ hence } \mathbb{E}[X] = 1/2$$
    $$\mathbb{E}[Z] = \mathbb{E}[N] \mathbb{E}[X] = 2 \cdot \frac{1}{2} = 1$$
\end{example}

\begin{remark}
    The Waldsche identity can be used to analyse the expected runtime of randomized algorithms. Where $N$ is the number of calls of the subroutine and $X$ the runtime of a single call. We can caluclate the total expected runtime using $Z$.
\end{remark}

\subsection{Distributions}

\begin{definition}
    Let $X$ be a random variable with values in $\{0,1\}$ and density
    $$P[X=1] = p \quad \text{and} \quad P[X=0] = 1-p$$
    Then $X$ follows a \textbf{Bernoulli distribution} with parameter $p$, denoted as $X \sim \text{Bernoulli}(p)$.
\end{definition}

\begin{example}
    Throwing a coin, indicatorvariable for heads.
\end{example}

\begin{remark}
    Observe that for a bernoulli-distribution we have:
    $$\mathbb{E}[X] = p \quad \text{and} \quad Var[X] = p(1-p)$$
\end{remark}

\begin{definition}
    Let $X$ be a random variable with values in $\{0,1,\dots,n\}$ and density
    $$P[X=k] = \binom{n}{k} p^k (1-p)^{n-k}$$
    Then $X$ follows a \textbf{Binomial distribution} with parameters $n$ and $p$, denoted as $X \sim \text{Bin}(n,p)$.
\end{definition}

\begin{example}
    Number of heads after throwing a coin $n$ times.
\end{example}

\begin{remark}
    Observe that for a binomial-distribution we have:
    $$\mathbb{E}[X] = np \quad \text{and} \quad Var[X] = np(1-p)$$
\end{remark}

\begin{definition}
    Let $X$ be a random variable with density
    $$P[X=k] = \frac{e^{-\lambda} \lambda^k}{k!}$$
    Then $X$ follows a \textbf{Poisson distribution} with parameter $\lambda$, denoted as $X \sim \text{Po}(\lambda)$.
\end{definition}

\begin{remark}
    Observe that for a poisson-distribution we have:
    $$\mathbb{E}[X] = \lambda \quad \text{and} \quad Var[X] = \lambda$$
\end{remark}

\begin{remark}
    This is often used to model events that do not happen often. For example the number of emails you get in a day. ($\lambda$ would be the average num of emails you get a day).
\end{remark}

\begin{definition}
    Let $X$ be a random variable with density
    $$P[X=k] = p(1-p)^{k-1}$$
    Then $X$ follows a \textbf{Geometric distribution} with parameter $p$, denoted as $X \sim \text{Geo}(p)$.
\end{definition}

\begin{example}
    Assume we throw a coin with $P[\text{head}] = p$ until we get n times heads. Let $X$ be the number of throws until we get n times heads.
    $$P[X=k] = \binom{k-1}{n-1} p^n (1-p)^{k-n}$$
\end{example}

\begin{remark}
    Observe that for a geometric-distribution we have:
    $$\mathbb{E}[X] = \frac{1}{p} \quad \text{and} \quad Var[X] = \frac{1-p}{p^2}$$
\end{remark}

\begin{definition}
    For $n \geq 2$ a random variable $X$ with density
    $$P[X=k] = \binom{k-1}{n-1} p^n (1-p)^{k-n}$$
    is called negative binomial distributed with order $n$.
\end{definition}

\begin{example}
    \textbf{Coupon Collector's Problem (Sammelbilder-Problem)}:
    There are $n$ different types of collectible cards. In every round we uniformly get one of them. Let $X$ be the number of rounds until we have all $n$ types of cards. What is $\mathbb{E}[X]$?
    
    We divide the process into phases: Phase $i$ denotes the steps from acquiring the $(i-1)$-th unique card (exclusive) to acquiring the $i$-th unique card (inclusive). Let $X_i$ be the number of rounds spent in phase $i$. Clearly, the total number of rounds is $X = \sum_{i=1}^n X_i$.
    
    In phase $i$, we already have $i-1$ unique cards. The probability of obtaining a new unique card in any round is:
    $$p_i = \frac{n - (i - 1)}{n} = \frac{n - i + 1}{n}$$
    Thus, $X_i$ is geometrically distributed with parameter $p_i$, meaning its expectation is:
    $$\mathbb{E}[X_i] = \frac{1}{p_i} = \frac{n}{n - i + 1}$$
    By the linearity of expectation, we have:
    $$\mathbb{E}[X] = \sum_{i=1}^n \mathbb{E}[X_i] = \sum_{i=1}^n \frac{n}{n - i + 1} = n \sum_{j=1}^n \frac{1}{j} = n H_n$$
    where $H_n$ is the $n$-th harmonic number. Since $H_n = \ln n + \mathcal{O}(1)$, we get:
    $$\mathbb{E}[X] = n \ln n + \mathcal{O}(n).$$
\end{example}

\subsection{Bounding Probability}

\begin{theorem}
    \textbf{Markov's Inequality}: Let $X$ be a \textbf{non-negative} random variable. For any $t \geq 0$:
    $$P[X \geq t] \leq \frac{\mathbb{E}[X]}{t}$$
    or equivalently:
    $$P[X \geq t \mathbb{E}[X]] \leq \frac{1}{t}$$
\end{theorem}

\begin{proof}
    $$\mathbb{E}[X] = \sum_{\alpha \in W_X} \alpha P[X=\alpha] = \sum_{\alpha \le t} \alpha P[X=\alpha] + \sum_{\alpha > t} \alpha P[X=\alpha]$$
    $$\geq \sum_{\alpha > t} \alpha P[X=\alpha] > \sum_{\alpha > t} t P[X=\alpha] = t P[X > t]$$
    $$\implies P[X > t] < \frac{\mathbb{E}[X]}{t}$$
\end{proof}

\begin{theorem}
    \textbf{Chebyshev's Inequality}: Let $X$ be a random variable. For any $t \geq 0$:
    $$P[|X-\mathbb{E}[X]| \geq t] \leq \frac{Var[X]}{t^2}$$
\end{theorem}

\begin{proof}
We want to show $P[|X-\mathbb{E}[X]| \geq t] \leq \frac{Var[X]}{t^2}$.
Let $Y = (X-\mathbb{E}[X])^2$. By Markov's inequality:
$$P[Y \geq t^2] \leq \frac{\mathbb{E}[Y]}{t^2}  \implies  P[(X-\mathbb{E}[X])^2 \geq t^2] \leq \frac{\mathbb{E}[(X-\mathbb{E}[X])^2]}{t^2}$$
$$\implies P[|X-\mathbb{E}[X]| \geq t] \leq \frac{Var[X]}{t^2}$$
\end{proof}

\begin{example}
    Let there be $n$ diffrent collectable items. We buy items $X$ times until we have a full set. We have already shown that $\mathbb{E}[X] = n H_n$. Now what is the variance of $X$?
    We know:
    $$X = \sum_{i=1}^n X_i \text{independent, geometric distributed with } p_i = \frac{n-i+1}{n}$$
    Hence
    $$Var[X_i] = \frac{1-p_i}{p_i^2} \leq \frac{1}{p_i^2} = \frac{n^2}{(n-i+1)^2}$$
    Further we show
    $$Var[X] = \sum_{i=1}^n Var[X_i] \leq \sum_{i=1}^n \frac{n^2}{(n-i+1)^2} = n^2 \sum_{i=1}^n \frac{1}{i^2} \leq n^2 \frac{\pi^2}{6}$$
    Using chebyshev's inequality we can bound the probability of $X$ being far from its mean. Let $t = n  \sqrt{\ln n}$:
    $$P[|X-\mathbb{E}[X]| \geq n \sqrt{\ln n}] \leq \frac{Var[X]}{t^2} \leq \frac{n^2 \frac{\pi^2}{6}}{n^2 \ln n} = \frac{\pi^2}{6 \ln n}$$
    
\end{example}

\begin{example}
    \textbf{Special case Binomial distribution: } Let $X \sim Bin(n, 1/2)$. Then $\mathbb{E}[X] = n/2$ and $Var[X] = n/4$.
    Using Markov's inequality:
    $$P[X \geq n] \leq \frac{\mathbb{E}[X]}{n} = \frac{n/2}{n} = 1/2$$
    Using Chebyshev's inequality:
    $$P[X \geq n] = P[|X-\mathbb{E}[X]| \geq n/2] \leq \frac{Var[X]}{(n/2)^2} = \frac{n/4}{n^2/4} = 1/n$$
\end{example}

\subsubsection{Chernoff Bounds}

\begin{remark}
    Chernoff Bounds can be used for binomial distribution $X \sim Bin(n, p)$.
\end{remark}

\begin{theorem}
    Let $X_1, \dots, X_n$ be independent Bernoulli random variables with $P[X_i=1] = p_i$. Let $X = \sum_{i=1}^n X_i$.
    $$P[X \geq (1+\delta)\mathbb{E}[X]] \leq e^{-\frac{1}{3}\delta^2\mathbb{E}[X]} \quad \forall 0 \leq \delta \leq 1$$
    $$P[X \leq (1-\delta)\mathbb{E}[X]] \leq e^{-\frac{1}{2}\delta^2\mathbb{E}[X]} \quad \forall 0 \leq \delta \leq 1$$
    $$P[X \geq t] \leq 2^{-t} \quad \forall t \geq 2 e \mathbb{E}[X]$$
\end{theorem}

\begin{proof}
    We prove the third inequality. We use Markov's inequality on $Y = 4^X$.
    We have 
    $$X \geq t \iff 4^X \geq 4^t$$
    Therfore it holds that
    $$P[X \geq t] = P[4^X \geq 4^t] \leq \frac{\mathbb{E}[4^X]}{4^t}$$
    Now what is $\mathbb{E}[4^X]$? Observe that as $X_1, X_2, \dots, X_n$ are independent, $4^{X_1}, 4^{X_2}, \dots, 4^{X_n}$ are also independent. Hence
    $$\mathbb{E}[4^X] = \mathbb{E}[4^{\sum X_i}] = \mathbb{E}[\prod_{i=1}^n 4^{X_i}] = \prod_{i=1}^n \mathbb{E}[4^{X_i}]$$
    Further
    $$\mathbb{E}[4^{X_i}] = \sum_{x=0}^n 4^x P[X_i=x] = (1-p_i) 4^0 + p_i 4^1 = 1 + 3p_i \leq e^{3p_i} \leq 2^t$$
    $$\implies \mathbb{E}[4^X] \leq \prod_{i=1}^n e^{3p_i} = e^{3\sum p_i} = e^{3\mathbb{E}[X]}$$
    $$\implies P[X \geq t] \leq \frac{\mathbb{E}[4^X]}{4^t} \leq \frac{2^t}{4^t} = \frac{1}{2^t}$$
\end{proof}

\begin{remark}
    To prove the other bounds we use the same technique on $Y = (1+\delta)^X$ and $Y = (1-\delta)^X$.
\end{remark}

\subsubsection{Target Shooting}

Given finite sets $S, U$ with $S \subseteq U$ we want to know $\frac{|S|}{|U|}$.
\begin{example}
    Let $U$ be the population of Switzerland and $S$ the set of people that have the flu.
\end{example}

We can't compute this directly, so we sample $k$ elements from $U$ and count how many of them are in $S$.

\begin{algorithm}
Choose u_1 ... u_n independently and uniformly at random from U
Return 1/n * |{u_i in S}| 
\end{algorithm}

Let $Y$ be the number of elements in $S$. We want to bound the approximation error $|Y - \frac{|S|}{|U|}|$. Observe that for every element $u_i$ we have an indipendent bernoulli variable $Y_i \in \{0,1\}$ such that $Y = \frac{1}{n} \sum_{i=1}^n Y_i$ and $P[Y_i=1] = \frac{|S|}{|U|}$.

\begin{theorem}
    Let $\delta, \epsilon \ge 0$ be given. If we choose $n \geq 3 \frac{|U|}{|S|} \frac{1}{\epsilon^2} \ln \frac{2}{\delta}$, then with probability at least $1-\delta$ we have that $Y$ is in 
    $$[(1-\epsilon)\frac{|S|}{|U|}, (1+\epsilon)\frac{|S|}{|U|}]$$
\end{theorem}

\begin{proof}
    We show that the probability of being outside that intervall is less than $\delta$.
    $$P\left[\left|Y - \frac{|S|}{|U|}\right| \geq \epsilon \frac{|S|}{|U|}\right] = P\left[\left|Y - \mathbb{E}[Y]\right| \geq \epsilon \mathbb{E}[Y]\right]$$
    $$ = P[|NY - \mathbb{E}[NY]| \geq n \epsilon \mathbb{E}[NY]]$$
    Now observe that $NY = Y_1 + \dots + Y_n$ where $Y_i$ are independent bernoulli variables with $P[Y_i=1] = \frac{|S|}{|U|}$. Hence we can use chernoff bounds.
    $$P[|NY - \mathbb{E}[NY]| \geq n \epsilon \mathbb{E}[NY]] \leq 2 e^{-\frac{1}{3} \epsilon^2 \mathbb{E}[NY]} =  2 e^{-\frac{1}{3} \epsilon^2 N \frac{|S|}{|U|}}$$
    With our assumption of $N \geq 3 \frac{|U|}{|S|} \frac{1}{\epsilon^2} \ln \frac{2}{\delta}$ we get exactly:
    $$P\left[\left|Y - \frac{|S|}{|U|}\right| \geq \epsilon \frac{|S|}{|U|}\right] \leq \delta$$
\end{proof}